% ── PREFACE ──────────────────────────────────────────────────────────────────
\chapter*{Preface}
\addcontentsline{toc}{chapter}{Preface}   % manually add to TOC (chapter* omits it)
\markboth{Preface}{Preface}               % set running header

\bigskip

\noindent\textit{This preface is a working placeholder. The final version will
be approximately 1{,}500 words, written in the first person. The four section
headings below mark the intended structure; the bracketed notes are drafting
prompts, not final text.}

\bigskip\bigskip

% ─────────────────────────────────────────────────────────────────────────────
\section*{How This Project Began}
% ─────────────────────────────────────────────────────────────────────────────

\noindent\textit{[Draft --- approx.\ 400 words. Personal narrative: the
single observation that motivated the book --- three independent traditions
arriving at the same logical structure. What it felt like to notice the
convergence. Why it demanded a book rather than just academic papers. Tone:
first-person, accessible; this is the reader's entry point.]}

\bigskip

% ─────────────────────────────────────────────────────────────────────────────
\section*{Who This Book Is Written For}
% ─────────────────────────────────────────────────────────────────────────────

\noindent\textit{[Draft --- approx.\ 300 words. Four audiences: (1)
scientists curious about what quantum mechanics implies for questions of
meaning, consciousness, and love; (2) philosophers working on the mind-body
problem or the philosophy of physics; (3) Bahá'í scholars and practitioners
interested in the two-wings principle; (4) general readers who suspect the
science-religion conflict is a false choice. Note which chapters each
audience may read selectively and which require prior chapters.]}

\bigskip

% ─────────────────────────────────────────────────────────────────────────────
\section*{How to Read This Book}
% ─────────────────────────────────────────────────────────────────────────────

\noindent\textit{[Draft --- approx.\ 400 words. Structure overview: Part I
introduces the three traditions; Part II builds the process-ontology bridge;
Part III documents the six convergences; Part IV develops the formal
mathematics and draws the implications. Each chapter ends with a Convergence
Moment box that restates the chapter's argument from all three traditions in
parallel. Parts I and II can be read selectively by specialists; Part III
depends on Part II. Physicists may enter at Chapter V; philosophers at
Chapter IV; Bahá'í readers at Chapter III.]}

\bigskip

% ─────────────────────────────────────────────────────────────────────────────
\section*{What This Book Is Not Claiming}
% ─────────────────────────────────────────────────────────────────────────────

\noindent\textit{[Draft --- approx.\ 400 words. Pre-empt the two main
misreadings. First: ``You are claiming the Bahá'í writings anticipated
quantum mechanics.'' No. The argument is structural convergence, not
prophetic foreknowledge. Three independent instruments registering the same
signal is not the same as one instrument predicting what the others would
find. Second: ``You are claiming love is just a quantum effect.'' No. The
argument runs the opposite direction: quantum entanglement is one formal
instance of what love, in its full ontological sense, has always named.
Physics does not reduce love; love names something that physics, at its
frontier, is finally precise enough to touch.]}

\bigskip

% ─────────────────────────────────────────────────────────────────────────────
\section*{Acknowledgements}
% ─────────────────────────────────────────────────────────────────────────────

\noindent\textit{[Draft --- to be written.]}

\vfill
